\documentclass[11pt,a4paper]{article}

\usepackage[utf8]{inputenc}
\usepackage[T1]{fontenc}
\usepackage{lmodern}
\usepackage[margin=2.5cm]{geometry}
\usepackage{graphicx}
\usepackage{booktabs}
\usepackage{amsmath}
\usepackage{microtype}
\usepackage{caption}
\usepackage[hidelinks]{hyperref}

% Allow long file names (typeset verbatim) to break at any character so they
% never run past the margin.
\urlstyle{tt}
\makeatletter
\g@addto@macro{\UrlBreaks}{\do\_\do\.\do\-\do\/\UrlOrds}
\makeatother
\newcommand{\fn}[1]{\nolinkurl{#1}}

\title{\textbf{A Historical Analysis of the Vienna Physics Curricula}\\[2mm]
\large Dating the Topics of the Bachelor's and Master's Programmes in Physics
(Version 2026)}
\author{Faculty of Physics, University of Vienna}
\date{\today}

\begin{document}
\maketitle

\begin{abstract}
\noindent
This report documents an analysis of the curricula of the Bachelor's and
Master's programmes in Physics at the University of Vienna (both Version 2026).
For each curriculum the individual scientific topics of all modules were
extracted from the official curriculum documents and decomposed into atomic
items. Each topic was then annotated with the year in which the corresponding
concept was first introduced or the corresponding discovery was made, together
with the responsible scientist(s) and a bibliographic source for the
attribution. The resulting datasets---$545$ topics for the Bachelor's programme
and $174$ for the Master's programme---were analysed statistically and
visualised as histograms of topics per decade. The Bachelor's curriculum spans
more than four centuries of physics and mathematics (median year $1905$),
whereas the Master's curriculum is sharply concentrated in the twentieth
century (median year $1936$), reflecting its role as a research-oriented
specialisation built on the broad foundations laid in the Bachelor's degree.
\end{abstract}

\section{Objective}
The goal of this analysis was to characterise the two Physics curricula not by
their administrative structure but by the \emph{historical depth} of their
scientific content. Concretely, for every topic taught we asked: \emph{when},
and \emph{by whom}, was this piece of physics or mathematics first established?
Aggregating the answers reveals how each programme is distributed across the
history of science and makes the pedagogical difference between the Bachelor's
and Master's degrees quantitatively visible.

\section{Data and Method}

\subsection{Source material}
The analysis is based on the two official curriculum documents,
\texttt{BA\_Physik.pdf} (Bachelor) and \texttt{MA\_Physics.pdf} (Master), both in
the Version~2026 draft. The text of every module description was extracted and
the enumerated content (the ``\textit{Die Inhalte umfassen}\,'' / possible-content
lists) was used as the basis for the topic catalogue. Learning-outcome and
competence statements were deliberately excluded, as they describe skills rather
than datable scientific concepts.

\subsection{Topic decomposition}
The curriculum texts frequently bundle several distinct concepts into a single
phrase. Every such phrase was decomposed into \emph{atomic} topics, each
corresponding to a single discovery or concept that can carry one year and one
attribution. For example, the Bachelor entry
``\textit{Gradient, Divergenz, Rotation}'' was split into three separate rows,
and the large elective modules \textit{Wahl-Fachbereiche} and
\textit{Wahllaborpraktikum} were unfolded into their constituent named methods
(e.g.\ FFT, quantum error correction, BCS superconductivity). The elective
subject areas were retained as an explicit grouping column so that the topics of
each field (Computational Physics, Quantum Information, Particle Physics, \dots)
remain identifiable.

\subsection{Annotation}
Each atomic topic received four annotations:
\begin{itemize}
  \item \textbf{Year} --- the year of the seminal publication that first
        introduced the concept or reported the discovery.
  \item \textbf{Scientist(s)} --- the originator(s); where a topic combines
        several contributions, the principal one is named with secondary
        contributors noted in parentheses.
  \item \textbf{Source} --- a bibliographic reference for the attribution,
        given as \emph{author, landmark work, year}. Publication titles are kept
        in their original language, which is standard scholarly practice.
  \item \textbf{Subject area} (electives only) --- the field heading under which
        the topic is offered.
\end{itemize}
Topics without a single identifiable originator---laboratory and presentation
skills, orientation courses, generic programming setup, the master's thesis and
defence, and similarly broad method headings---were marked ``---'' and excluded
from the statistical analysis.

\subsection{Deliverables and reproducibility}
For each curriculum the catalogue was produced as a spreadsheet in both German
and English (identical content, translated headers, module names, subject areas
and topics). The decade histograms are generated directly from the spreadsheets
by two short, self-contained Python scripts
(\texttt{make\_histograms.py} and \texttt{make\_combined\_histogram.py}), so the
figures can be regenerated at any time. The complete list of files is given in
Section~\ref{sec:files}.

\subsection{Caveats}
The attributions are, by their nature, scholarly judgements. For foundational
results with multiple independent originators, and especially for the broad
elective and core subject areas that represent entire research fields, a single
\emph{representative} seminal contribution was chosen. The source column
documents the work that established each concept to the level of
\emph{author--title--year}; exact volume and page numbers were intentionally
omitted to avoid spurious precision. Individual attributions can therefore be
revised without affecting the overall picture.

\section{The Topic Catalogues}

\subsection{Bachelor's programme}
The Bachelor catalogue contains $545$ atomic topics across $25$ modules, of
which $476$ carry a year and $69$ are non-attributable. The two large elective
modules dominate in size: \textit{Wahl-Fachbereiche} contributes $149$ topics
across eleven subject areas and \textit{Wahllaborpraktikum} a further $62$ across
eight. The core lecture modules (Experimental Physics, Theoretical Physics,
Analysis, Linear Algebra, \dots) contribute the remainder.

\subsection{Master's programme}
The Master catalogue contains $174$ atomic topics across $8$ modules, of which
$140$ carry a year and $34$ are non-attributable. Here the scientific content is
concentrated almost entirely in a single module: \textit{Core Topics} provides
$150$ topics, organised into $15$ subject areas
(Atmospheric Aerosol Physics, Biophysics, Computational Physics, Condensed
Matter Physics, Electronic Structure, Experimental Quantum Physics, General
Relativity and Cosmology, Quantum Field Theory, Phase Transitions and Critical
Phenomena, Physics of Nuclei and Isotopes, Non-equilibrium Statistical Physics,
Scattering/Microscopy/Spectroscopy, Soft Matter Physics, Theoretical Particle
Physics, and Theory in Quantum Optics and Quantum Information). The remaining
modules (Scientific Research Methods, Advanced Research Topics, Electives,
Proposal/Master's Seminar, Thesis and Defensio) are research- and
skills-oriented and contribute few datable topics.

\section{Results}

\subsection{Summary statistics}
Table~\ref{tab:summary} summarises the distribution of topic years for the two
programmes. For the Bachelor's programme the single ancient entry
(Euclidean geometry, $\approx 300$\,BC) is treated as an outlier and excluded
from the summary figures; it is the only pre-Renaissance topic and, when
included, shifts the mean down to $1877$ and inflates the standard deviation
from $82$ to $129$ years.

\begin{table}[h]
\centering
\caption{Summary statistics of the topic years (dated topics only).
Bachelor figures exclude the single Euclidean-geometry outlier.}
\label{tab:summary}
\begin{tabular}{lcc}
\toprule
Quantity & Bachelor & Master \\
\midrule
Dated topics $n$        & $475$              & $140$ \\
Median year             & $1905$             & $1936$ \\
Mean year               & $1882$             & $1936$ \\
Std.\ deviation (years) & $82$               & $35$ \\
Earliest topic          & $1572$ (Bombelli)  & $1828$ (Green) \\
Latest topic            & $2004$ (graphene)  & $2001$ (KLM) \\
Modal decade            & 1920s ($51$)       & 1930s ($22$) \\
\bottomrule
\end{tabular}
\end{table}

\subsection{Bachelor's programme}
The distribution of Bachelor topics by decade is shown in
Figure~\ref{fig:ba}. The curriculum covers the entire history of modern physics
and its mathematical underpinnings: a sparse seventeenth-- and
eighteenth--century tail (Newton, Leibniz, Euler, Lagrange), a pronounced rise
through the nineteenth century (classical electrodynamics, thermodynamics, and
the rigorous foundations of analysis), and a dominant early--twentieth--century
peak in the \textbf{1920s} ($51$ topics) driven by the quantum--mechanical
revolution. Content continues at a substantial level up to the present day
through quantum information, condensed-matter and computational topics.
The oldest non-ancient topic is the introduction of complex numbers
(Bombelli, $1572$); the most recent is graphene (Geim and Novoselov, $2004$).

\begin{figure}[h]
\centering
\includegraphics[width=0.92\textwidth]{BA_Physics_Topics_Histogram.png}
\caption{Number of Bachelor topics per decade. Vertical dashed lines mark
centuries. The distribution peaks in the 1920s and has a median year of $1905$.}
\label{fig:ba}
\end{figure}

\subsection{Master's programme}
The corresponding distribution for the Master's programme is shown in
Figure~\ref{fig:ma}. It is markedly more modern and far more tightly
concentrated: there is essentially no content before $\sim 1820$, and the mass
of the distribution lies in the twentieth century with a peak in the
\textbf{1930s} ($22$ topics). The median year is $1936$ and the standard
deviation is only $35$ years---less than half that of the Bachelor's
programme. The earliest topic is Green's functions (Green, $1828$); the most
recent is photonic quantum information processing
(Knill, Laflamme and Milburn, $2001$).

\begin{figure}[h]
\centering
\includegraphics[width=0.92\textwidth]{MA_Physics_Topics_Histogram.png}
\caption{Number of Master topics per decade. The distribution is concentrated in
the twentieth century with a median year of $1936$.}
\label{fig:ma}
\end{figure}

\subsection{Comparison}
Figure~\ref{fig:bama} overlays the two distributions on a common axis. The
contrast captures the pedagogical logic of the two-cycle degree structure: the
Master's distribution (orange) sits almost entirely beneath the modern end of
the Bachelor's distribution (blue). The Bachelor's programme builds the
historical and mathematical \emph{foundations} of physics across four centuries,
while the Master's programme assumes those foundations and specialises in
twentieth- and twenty-first-century \emph{research} physics. Numerically, the
centre of gravity moves forward by about three decades, from a median of $1905$
to $1936$.

\begin{figure}[h]
\centering
\includegraphics[width=0.98\textwidth]{BA_MA_Topics_Histogram.png}
\caption{Combined histogram of Bachelor (blue) and Master (orange) topics by
decade. The Master's content is confined to the modern end of the Bachelor's
range. Counts are raw; the Master's programme has fewer topics in total
($140$ vs.\ $475$ dated topics).}
\label{fig:bama}
\end{figure}

\section{Discussion}
The quantitative results confirm and sharpen an intuitive expectation. A
Bachelor's curriculum in physics must, of necessity, recapitulate the historical
development of the field: one cannot teach quantum mechanics without first
teaching classical mechanics, calculus, and electrodynamics. This produces a
broad distribution with a long tail reaching back to the seventeenth century
(and, for Euclidean geometry, to antiquity). A research-oriented Master's
curriculum, by contrast, can take this entire edifice as given and devote itself
to the modern research frontier; its topic distribution is therefore both later
in time and far narrower. The roughly $30$-year forward shift of the median, and
the halving of the temporal spread, are compact quantitative signatures of this
transition from \emph{foundational breadth} to \emph{research-level depth}.

\section{Limitations}

\paragraph{Earliest-date bias.}
The most important limitation concerns the meaning of the assigned year. By
construction, each topic is dated by the \emph{earliest} possible date---the
year in which the concept was first introduced or the founding discovery was
made---whereas a course actually treats the topic together with all of its
subsequent developments. The year therefore marks the \emph{birth} of a subject,
not the full historical range of the material taught under that heading.

A clear example is the topic \textit{Molecular dynamics} in the Computational
Physics course, which is dated to $1957$, the year of the first hard-sphere
molecular-dynamics simulation by Alder and Wainwright. The course, however, does
not stop there: it also covers much more recent methodological developments such
as thermostat and barostat algorithms for sampling different ensembles
(1980s), constraint algorithms, and symplectic / time-reversible integration
schemes, many of which date from the 1980s and later. All of this modern content
is nevertheless represented in the dataset by the single founding year $1957$.

The same effect applies throughout: a topic such as ``density functional
theory'' (dated $1965$) is taught together with exchange--correlation functionals
and methods developed decades later, and ``general relativity'' (dated $1915$)
is accompanied by modern results on gravitational waves and numerical relativity.
As a consequence, the histograms \emph{systematically shift weight towards
earlier decades} and should be read as the distribution of the \emph{conceptual
origins} of the curriculum, not as the distribution of the actual age of the
material presented in the lecture halls. The true content is, on average,
more modern than the founding-year distributions suggest. A more faithful but
considerably more laborious alternative would be to assign each topic a range of
years, or to record separately the principal later milestones; this was beyond
the scope of the present analysis.

\paragraph{Further limitations.}
Three additional points should be kept in mind. First, the year and scientist
assignments are scholarly judgements; for fields rather than single discoveries
a representative seminal work was chosen. Second, the histograms weight every
atomic topic equally, irrespective of the teaching time (ECTS) devoted to it, so
the analysis measures the \emph{conceptual} rather than the \emph{temporal}
balance of the curriculum. Third, the decomposition of compound topics into
atomic items necessarily involves some discretion, which affects the absolute
counts (but not the qualitative shape of the distributions).

\section{Files Produced}
\label{sec:files}
The analysis produced the following deliverables, all in the project directory:
\begin{itemize}
  \item \textbf{Spreadsheets (Bachelor):}
        \fn{BA_Physik_Module_Themen.xlsx} (German),
        \fn{BA_Physics_Modules_Topics.xlsx} (English).
  \item \textbf{Spreadsheets (Master):}
        \fn{MA_Physik_Module_Themen.xlsx} (German),
        \fn{MA_Physics_Modules_Topics.xlsx} (English).
  \item \textbf{Histograms (Bachelor):}
        \fn{BA_Physik_Themen_Histogramm.png} (German),
        \fn{BA_Physics_Topics_Histogram.png} (English).
  \item \textbf{Histograms (Master):}
        \fn{MA_Physik_Themen_Histogramm.png} (German),
        \fn{MA_Physics_Topics_Histogram.png} (English).
  \item \textbf{Combined histograms:}
        \fn{BA_MA_Themen_Histogramm.png} (German),
        \fn{BA_MA_Topics_Histogram.png} (English).
  \item \textbf{Scripts:}
        \fn{make_histograms.py},
        \fn{make_combined_histogram.py}.
\end{itemize}
Each spreadsheet contains a \textit{Topics by Module} sheet (one row per atomic
topic, with module, subject area, topic, year, scientist and source) and a
\textit{Module Overview} sheet (topic counts per module).

\end{document}
